\documentclass[UTF8,11pt,openany,fontset=fandol]{ctexbook}
\usepackage[paperwidth=170mm,paperheight=240mm,inner=22mm,outer=19mm,top=23mm,bottom=24mm,headheight=15pt]{geometry}
\usepackage{amsmath,amssymb,amsthm,mathtools}
\usepackage{booktabs,longtable,array,tabularx}
\usepackage{graphicx,xcolor,tikz}
\usetikzlibrary{arrows.meta,positioning,calc,fit,backgrounds}
\usepackage{enumitem,fancyhdr,microtype}
\usepackage[numbers,sort&compress]{natbib}
\usepackage{hyperref}
\definecolor{ink}{HTML}{17334B}
\definecolor{accent}{HTML}{A45335}
\definecolor{pale}{HTML}{E9EFF3}
\hypersetup{unicode=true,colorlinks=true,linkcolor=ink,citecolor=accent,urlcolor=ink,pdftitle={思维论},pdfsubject={诊断组织的有限数学：共同错误、可辨别性与通信},pdfauthor={AI员工项目研究书稿}}
\urlstyle{same}
\setlength{\emergencystretch}{2em}
\setlength{\parskip}{0.25em}
\setlength{\tabcolsep}{5pt}
\linespread{1.12}
\raggedbottom
\setcounter{tocdepth}{1}
\setcounter{secnumdepth}{2}
\setlist{nosep,topsep=4pt,leftmargin=2em}
\pagestyle{fancy}
\fancyhf{}
\fancyhead[LE,RO]{\small\thepage}
\fancyhead[RE]{\small 思维论}
\fancyhead[LO]{\small\nouppercase{\leftmark}}
\renewcommand{\headrulewidth}{0.3pt}
\ctexset{chapter={format=\Large\bfseries\color{ink},beforeskip=12pt,afterskip=22pt},section={format=\large\bfseries\color{ink}}}
\theoremstyle{definition}
\newtheorem{definition}{定义}[chapter]
\newtheorem{example}[definition]{例}
\newtheorem{counterexample}[definition]{反例}
\newtheorem{exercise}{练习}[chapter]
\theoremstyle{plain}
\newtheorem{theorem}[definition]{定理}
\newtheorem{proposition}[definition]{命题}
\newtheorem{lemma}[definition]{引理}
\newtheorem{corollary}[definition]{推论}
\theoremstyle{remark}
\newtheorem{remark}[definition]{说明}
\renewcommand{\proofname}{证明}
\newcommand{\VV}{\mathcal V}
\newcommand{\HH}{\mathcal H}
\newcommand{\LL}{\mathcal L}
\newcommand{\XX}{\mathcal X}
\newcommand{\II}{\mathbf I}
\newcommand{\PP}{\mathbb P}
\newcommand{\EE}{\mathbb E}
\newcommand{\NN}{\mathbb N}
\newcommand{\RR}{\mathbb R}
\newcommand{\concat}{\mathbin{\Vert}}
\newcommand{\wait}{\mathsf{wait}}
\newcommand{\StopMode}{\mathsf{stop}}
\newcommand{\TV}{\operatorname{TV}}
\newcommand{\Var}{\operatorname{Var}}
\newcommand{\Cov}{\operatorname{Cov}}
\newcommand{\KL}{D_{\mathrm{KL}}}
\newcommand{\bits}{\operatorname{bits}}
\newcommand{\cost}{\operatorname{cost}}
\newcommand{\obs}{\operatorname{obs}}
\newcommand{\sem}[1]{[\![#1]\!]}
\newcommand{\chapterguide}[1]{\begin{quote}\small\color{ink}#1\end{quote}}
\DeclareMathOperator{\Risk}{Risk}
\DeclareMathOperator{\Reg}{Reg}
\DeclareMathOperator{\Bin}{Bin}
\DeclareMathOperator{\Bern}{Bern}
\newcommand{\calP}{\mathfrak P}
\newcommand{\calD}{\mathcal D}
\newcommand{\calA}{\mathcal A}
\newcommand{\calQ}{\mathcal Q}
\newcommand{\kap}{\kappa_{\mathrm H}}
\newcommand{\kapK}{\kappa_{\mathrm K}}
\newcommand{\Aff}{\operatorname{Aff}}
\newcommand{\Alt}{\operatorname{Alt}}
\newcommand{\kl}{\operatorname{kl}}
\newcolumntype{Y}{>{\raggedright\arraybackslash}X}
\tikzset{fragment/.style={circle,draw=ink,fill=pale,minimum size=7mm,inner sep=1pt},operation/.style={rectangle,draw=accent,fill=white,minimum width=9mm,minimum height=6mm},unit/.style={rounded corners,draw=ink,fill=pale,minimum width=22mm,minimum height=11mm,align=center},edge/.style={-{Stealth[length=2mm]},semithick,draw=ink}}

\title{思维论}
\begin{document}
\frontmatter
\begin{titlepage}
\thispagestyle{empty}
\vspace*{30mm}
{\fontsize{40}{50}\selectfont\bfseries\color{ink}思维论\par}
\vspace{15mm}
\begin{tikzpicture}[node distance=17mm]
\node[fragment] (input) {$v_0$};
\node[operation,above right=7mm and 14mm of input] (first) {$F_1$};
\node[operation,below right=7mm and 14mm of input] (second) {$F_2$};
\node[fragment,right=14mm of first] (left) {$v_1$};
\node[fragment,right=14mm of second] (right) {$v_2$};
\node[operation,below right=7mm and 14mm of left] (combine) {$F_3$};
\node[fragment,right=of combine] (output) {$v_3$};
\draw[edge] (input)--(first); \draw[edge] (first)--(left);
\draw[edge] (input)--(second); \draw[edge] (second)--(right);
\draw[edge] (left)--(combine); \draw[edge] (right)--(combine);
\draw[edge] (combine)--(output);
\end{tikzpicture}
\vfill
{\large 诊断组织的有限数学\par}
\vspace{3mm}
{共同错误、可辨别性与通信\par}
\vspace{5mm}
{A Theory of Organized Reasoning\par}
\end{titlepage}

\chapter*{前言}
\addcontentsline{toc}{chapter}{前言}
几个推理过程可能反复确认同一个答案，却始终没有排除造成这个答案的共同误解。可靠性取决于它们能够取得哪些区别、怎样分配取得区别的机会，以及这些区别能否到达真正需要它们的过程。增加重复次数只有在这些条件允许时才有效。

一个具体问题贯穿全书：\emph{当基本操作可能受到共同错误影响，控制者只能通过这些操作取得证据时，什么样的组织能够把错误概率持续压低，需要付出多少调用与通信？}

研究从有限的未知情境出发。情境同时规定真实目标与隐藏的失效因素。每个基本操作在不同情境下产生不同的片段分布；组织可以根据已经得到的片段决定下一步。这允许把错误相关性明确写入模型，并把所有自适应组织放进同一个比较范围。

第一条结果刻画永久障碍：两个需要不同答案的情境，如果在所有允许操作下都产生相同分布，任何组织都无法将它们分开。第二条结果给出建设性方向：用一份诊断分配覆盖所有目标冲突对，可以得到统一的可靠性上界。第三条结果利用自适应实验的相对熵账本，约束整个组织类别。在固定有限情境数和有界似然比的范围内，诊断覆盖率由此同时控制充分预算与必要预算的数量级。

带共同偏差的二元模型把这些结果变成可计算的例子。额外操作是否有用，由它能否打破混淆类决定；中间表示能否构造，可以通过有限域线性关系判断。通信章研究哪些消息保存诊断能力，哪些压缩会把可解任务重新变成不可解任务，以及何时已有证据足以结束交流。

统计实验、主动检验和有噪计算提供了数学基础。经典工具在使用处归属；本书展开有限模型中的连接、推导和可复算实例。全书的统一性由几条相互调用的结果承担：不可辨别类、诊断覆盖率、自适应下界，以及保存这些量的通信与表示变换。

\section*{阅读路线}
第1--2章建立对象与永久障碍；第3--5章给出覆盖率、上下界和主定理；第6章分析共同偏差与校准；第7章独立研究通信；第8章讨论表示、计算和模型误差；第9--10章据此解释认知现象并提出未解决的问题。附录提供基础工具、练习解答和核验说明。

\chapter*{符号与约定}
\addcontentsline{toc}{chapter}{符号与约定}
\begin{longtable}{p{0.23\textwidth}p{0.70\textwidth}}
\toprule
符号 & 含义\\\midrule
$v,\II_e,e$ & 思维片段、有序输入和一次推理操作。\\
$\LL_t,\Pi$ & 操作记录，以及决定后续操作的推理链模式。\\
$\theta\in\Theta$ & 未知情境，可以包含答案和共同错误因素。\\
$g(\theta)$ & 该情境下任务要求的答案。\\
$a,K_a(v\mid\theta)$ & 基本实验，以及它产生片段的概率。\\
$d,m$ & 情境数，以及任一情境最多有多少目标不同的备选情境。\\
$h_a(\theta,\theta')$ & 一个实验对两个情境的Hellinger区别量。\\
$\kap,\kapK$ & 全局诊断覆盖率与局部自适应KL参数。\\
$N_a(n),T_n$ & 实验调用次数与完整运行记录。\\
$n^*(\delta)$ & 最坏情境错误率不超过$\delta$的最少调用数。\\
$\mu(\calA,l)$ & 二元线性模型中的结构覆盖率。\\
\bottomrule
\end{longtable}
除特别说明，集合有限，$d\ge2$，$g$至少取两个值，$0<\delta<1/2$。$\ln$和$\KL$使用自然单位；$\log_2$、$H_2$和$I_2$使用bit。$\PP_\theta$表示固定真实情境后的运行概率。节点直接表示片段，不另设一套内容变量。目标可以由外部给出，每次操作无需重新生成目标。

\tableofcontents
\listoffigures
\mainmatter
\part{对象与障碍}
\chapter{有限诊断中的思维过程}
\label{ch:model}
\chapterguide{先固定系统能够观察什么，再比较它怎样组织观察。未知情境的真值不进入控制者输入。}

\section{一个持续的共同错误}
设目标为事实$Y\in\{0,1\}$，某条推理规则受到隐藏偏差$B\in\{0,1\}$影响。重复使用它可能稳定得到$Y\oplus B$。多数表决可以压低独立噪声，却无法决定稳定答案来自真实事实还是共同偏差。

每次调用的扰动会重新发生；偏差$B$则贯穿整段运行。用固定但未知的情境$\theta$同时记录真实目标与持续因素。不同成员即使使用独立随机种子，也可能共同受到$B$影响。

\section{片段、操作与推理链模式}
片段$v$可以是一项观察、一个候选解释或已经组合好的记录。一次推理写作
\[
 \II_e\xrightarrow{F_a}v,\qquad \II_e=[v_{i_1},\ldots,v_{i_r}].
\]
有序输入有向超边代表一次操作。两次独立调用分别产生两个节点，同一节点可以被多个操作读取。

推理链$\LL_t=[e_1,\ldots,e_t]$包含操作及其输入输出。推理链模式$\Pi$依据这条记录和外部到达的可用材料，选择下一次绑定了输入的操作、等待或停止。执行模式并持续保有自身历史与可见材料的运行实例称为Agent。模式可以由另一Agent提供；生成它所用的操作也进入总成本。

片段层自回归模式规定下一输入追加刚得到的片段；其他模式可以重新选择和排列片段。单次$F_a$内部是否逐token生成，属于另一层描述。

多Agent组织的成员根据各自局部历史行动。证明下界时，可以让更强的比较者看到全部记录；构造分布式算法时，必须说明记录通过什么消息到达所需位置。

\section{受控的有限实验族}
\begin{definition}[诊断实例]
实例由有限集合$\Theta,\calA,\VV$、目标函数$g:\Theta\to\mathcal Y$和概率表$K_a(v\mid\theta)$组成。第$t$次调用满足
\begin{equation}
 \PP_\theta(V_t=v\mid T_{t-1},A_t=a)=K_a(v\mid\theta),
 \quad T_t=[(A_1,V_1),\ldots,(A_t,V_t)].
 \label{eq:controlled}
\end{equation}
控制者知道概率表与目标函数，未知真实$\theta$，只能通过调用取得有关$\theta$的新信息。
\end{definition}

给定$\theta$与所选实验，下一观察不再依赖更早记录；选择哪个实验仍可依赖全部历史。持续错误通过$\theta$进入各次调用。对未知情境混合后，观察一般相关。

实验$a$是一种允许的查询或处理方式，有限输出$v$就是形成的片段。这一专门化模型研究证据取得与组织。历史会改变基本操作分布的系统，需要增添显式状态；第\ref{ch:frontiers}章讨论新的困难。

\section{完整组织类别与成本}
允许所有依据历史选择实验的随机策略，其内部随机性与真实情境独立。结束时输出$\widehat Y$。定义
\[
 R_n^*=\inf_{\Pi:\,N\le n\ \mathrm{a.s.}}
 \max_\theta\PP_\theta^\Pi(\widehat Y\ne g(\theta)),
 \qquad n^*(\delta)=\min\{n:R_n^*\le\delta\}.
\]
集合为空时$n^*(\delta)=\infty$。停止后用不产生信息的空操作补齐记录，便能同时分析提前停止策略。

主定理计数等成本基本调用。控制者的查表、计数、似然计算和通信另行报告。下界允许控制者使用任意计算，因此也约束受限控制者；上界给出实现所需运算。若调用成本为$c_a$，应增加$\sum_ac_aN_a\le B$约束，并用同样权重改写覆盖规划。

统计定理使用已知概率核。实现部分以显式有理概率表为输入，对平方根、对数和分配作有误差控制的近似，再核对实际覆盖下界。精度成本另计，不把任意实数视为免费可读的无限信息。第\ref{ch:representation}章分析这类数值误差如何进入风险。

\section{一个可以核对的运行}
取四个情境$(Y,B)\in\{0,1\}^2$，两个实验分别观察$Y\oplus B$与$B$，各以概率$\eta<1/2$翻转。策略先重复第一种实验，再决定是否校准第二种。最终输出从观察中得到，控制者不能直接读取$Y$或$B$。

这里没有预设永远可靠的检查器。操作能否区分需要不同答案的情境，将从概率表推出。
\begin{exercise}把全部随机性写入一个预先采样的变量，为什么不意味着控制者能够读取这个变量的真实值？\end{exercise}
\begin{exercise}构造两个拥有相同最后片段、但适合选择不同下一操作的历史，解释模式为何需要使用运行记录。\end{exercise}

\chapter{共同错误的不可辨别类}
\label{ch:identifiability}
\chapterguide{永久盲点由实验分布决定。调度与协作不能消除全部允许操作共同保留的混淆。}

\section{把全部操作一起比较}
定义等价关系
\[
 \theta\sim\theta'\quad\Longleftrightarrow\quad
 K_a(\cdot\mid\theta)=K_a(\cdot\mid\theta')\quad\text{对每个 }a.
\]
它把$\Theta$划分为不可辨别类。类内情境的解释可以不同，但全部允许观察的分布相同。

\begin{theorem}[任意组织的不变盲点]
\label{thm:blind}
若$\theta\sim\theta'$，任意有限预算自适应组织在两种情境下产生相同的完整记录分布。消息、压缩和最终输出也具有相同分布。
\end{theorem}
\begin{proof}
初始材料和策略随机性的分布相同。假设此前记录的分布相同，在同一记录上下一实验的选择概率就相同；给定实验，输出核又相同。因此扩展后记录的分布相同。归纳完成。对相同分布进行同样的后处理仍得到相同分布。
\end{proof}

\begin{theorem}[不可辨别类的精确渐近误差]
\label{thm:floor}
给定先验$p_\theta$，最优Bayes错误率随预算趋于无穷时的极限为
\[
 R_\infty(p)=\sum_{C\in\Theta/{\sim}}
 \left[p(C)-\max_y\sum_{\theta\in C:\,g(\theta)=y}p_\theta\right].
\]
\end{theorem}
\begin{proof}
即使额外告知类$C$，类内情境的记录似然仍相同，后验比保持先验比。输出类内先验质量最大的目标，得到所写下界。

为达到它，轮流重复每种实验。各经验分布几乎必然收敛到真实分布。不同类的完整概率表不同，且候选有限；不同表之间的最小正距离为正。用经验表最近的类作估计，误分到其他类的概率趋于零。然后输出类内先验最优目标，由损失有界得到极限。
\end{proof}

\begin{corollary}[目标可辨别性的充要条件]
\label{cor:identifiability}
对每个$\delta>0$都存在有限预算使最坏错误率不超过$\delta$，当且仅当$g$在每个不可辨别类上为常数。
\end{corollary}
\begin{proof}
若某类含不同目标，取两个冲突情境的等先验，错误率至少为$1/2$。反之，前述类识别过程在每个情境下错误趋于零；有限多个情境的最大错误也趋于零。
\end{proof}

\section{重复确认的极限}
只观察$Y\oplus B$，四个情境先验均匀。类为$\{(0,0),(1,1)\}$与$\{(0,1),(1,0)\}$，每类的两种目标各占一半，故$R_\infty=1/2$。重复可以可靠识别$Y\oplus B$；将其错误地解释成对$Y$的证据，会使系统对未被识别的目标越来越自信。

\begin{figure}[ht]
\centering
\begin{tikzpicture}
\node[unit] (first) {$(0,0)$\quad$(1,1)$\\共同观察中心：0};
\node[unit,right=12mm of first] (second) {$(0,1)$\quad$(1,0)$\\共同观察中心：1};
\node[below=8mm of first,align=center,font=\small] {目标分别为0和1\\重复不能分开};
\node[below=8mm of second,align=center,font=\small] {目标分别为0和1\\重复不能分开};
\end{tikzpicture}
\caption{未校准的共同偏差形成两个不可辨别类。括号依次记录目标与偏差。}
\end{figure}

\section{关系与边界}
比较观察机制允许哪些决策，是统计实验比较的基本路线\citep{blackwell1953}。这里先处理完全等同的实验分布，后续再衡量区别的强弱。条件可从给定概率表有限检查；取得或识别这份概率表本身仍需要资源。

\begin{exercise}三个类的先验质量为$0.4,0.3,0.3$，类内目标0的条件概率为$1,1/2,1/3$。计算渐近最优错误率。\end{exercise}
\begin{exercise}扩大允许实验集合后，不可辨别类如何变化？证明渐近最优错误率不会增大。\end{exercise}

\part{组织能力的上下界}
\chapter{诊断覆盖率与可靠构造}
\label{ch:coverage}
\chapterguide{覆盖规划分配取得不同证据的机会，同时照顾所有需要不同答案的情境对。}

\section{一次观察提供多少区别}
对两个分布定义
\[
 \Aff(P,Q)=\sum_v\sqrt{P(v)Q(v)},\qquad h(P,Q)=1-\Aff(P,Q).
\]
Cauchy--Schwarz给出$0\le h\le1$，$h=0$恰当$P=Q$；又有$\sum_v(\sqrt P-\sqrt Q)^2=2h$。记$h_a(\theta,\theta')=h(K_a(\cdot\mid\theta),K_a(\cdot\mid\theta'))$。只要求区分目标不同的情境，令
\[
\begin{aligned}
 \mathcal E&=\{\{\theta,\theta'\}:g(\theta)\ne g(\theta')\},\\
 \Alt(\theta)&=\{\theta':g(\theta')\ne g(\theta)\},\\
 m&=\max_\theta|\Alt(\theta)|.
\end{aligned}
\]
\begin{definition}[诊断覆盖率]
设$\Delta(\calA)$为实验分配概率集合，定义
\begin{equation}
 \kap=\max_{w\in\Delta(\calA)}\min_{\{\theta,\theta'\}\in\mathcal E}
                  \sum_aw_ah_a(\theta,\theta').
 \label{eq:coverage-rate}
\end{equation}
\end{definition}
$w$规定每次调用各种实验的频率，内层最小值关注最难照顾的目标冲突对。它由基本概率表和目标定义，无需先知道最优组织。

\begin{proposition}[正覆盖率的精确条件]
\label{prop:positive}
$\kap>0$当且仅当$g$在不可辨别类上为常数。
\end{proposition}
\begin{proof}
目标冲突的不可辨别对具有全零区别量。反之，每个冲突对至少有一个正区别量；对所有实验用严格正的权重，每对加权和都正，有限多对的最小值仍正。
\end{proof}

\section{覆盖规划与反方证据}
最大化$u$，约束$w_a\ge0$、$\sum_aw_a=1$、每个冲突对满足$\sum_aw_ah_a\ge u$，便得到一份线性规划。其对偶形式为
\begin{equation}
 \kap=\min_{\lambda\in\Delta(\mathcal E)}\max_a
        \sum_{\{\theta,\theta'\}\in\mathcal E}\lambda_{\theta\theta'}h_a(\theta,\theta').
 \label{eq:dual}
\end{equation}
有限线性规划对偶在附录\ref{app:tools}说明。给出$w$可以证明覆盖率的下界；给出冲突对分配$\lambda$可以证明上界。二者相等时，最优值获得双向证据。

两种实验若分别只区分两类冲突，且正区别量都为$h_0$，则$w=(1/2,1/2)$给出$h_0/2$；反方等概率选择这两类冲突，也给出上界$h_0/2$。增加某一实验的比例会让另一类冲突成为瓶颈。

\section{一个无需知道真实情境的组织}
每步独立按固定$w$选择实验，保存$(a,v)$。计算
\[
 L_n(\vartheta)=\prod_{t=1}^nK_{A_t}(V_t\mid\vartheta),
\]
选择似然最大情境$\widehat\theta$，输出$g(\widehat\theta)$，平局按固定顺序处理。实验选择概率与$\vartheta$无关，在似然比较中消去。

\begin{theorem}[覆盖率的可靠性保证]
\label{thm:upper}
若$w$使每个冲突对的平均区别量至少为$u$，则
\[
 \max_\theta\PP_\theta(g(\widehat\theta)\ne g(\theta))
 \le\min\{1,m(1-u)^n\}\le\min\{1,me^{-nu}\}.
\]
$\kap>0$时，$n=\lceil\ln(m/\delta)/\kap\rceil$足够。
\end{theorem}
\begin{proof}
单步联合分布$P_\theta(a,v)=w_aK_a(v\mid\theta)$。两个情境的亲和度为$1-\sum_aw_ah_a\le1-u$，独立乘积的亲和度相乘。事件$L_n(\theta')\ge L_n(\theta)$上，真实记录概率不超过两个记录概率的几何平均，因此该事件概率至多$(1-u)^n$，零概率存在时也成立。输出错误必有一个$\theta'\in\Alt(\theta)$的似然不小于真值。对至多$m$个情境相加即得。
\end{proof}

\section{实现与近似}
分配从概率表求得，观测前无需知道真实情境。概率表尺寸为$|\calA|d|\VV|$；规划有$|\calA|+1$个变量与至多$d(d-1)/2$个冲突约束。在线对数似然更新每次需要$O(d)$次表读取和加法。

若只算出经核对的分配，覆盖率$u\ge\kap-\varepsilon>0$，使用$u$即可，无需把数值最优误认成精确最优。表的生成、规划求解与有限精度计算应与调用数一起报告。
\begin{exercise}两个冲突对在两个实验上的区别量为$(a,0)$与$(0,b)$，$a,b>0$。求最优分配与覆盖率。\end{exercise}
\begin{exercise}目标为$g(\theta)$时，为何可以用完整情境似然比较证明风险上界？这种分析可能多付出什么因子？\end{exercise}

\chapter{任意自适应组织的证据预算}
\label{ch:lower}
\chapterguide{下界比较整个允许的组织类别。自适应选择能够重新分配证据，不能免除区分冲突情境所需的证据量。}

\section{记录的似然比}
本章假设所有$K_a(v\mid\theta)>0$。这排除一次观察即绝对确定情境的特殊情形，使相对熵有限。令$P_\theta^\Pi$为组织完整记录的分布。

策略在历史$h$上选择实验的概率为$\pi_t(a\mid h)$。对同一条记录，这个概率在两个待比较情境下相同；虽然真实情境会改变历史的出现频率，它不直接进入策略函数。因此记录似然比为
\[
 \frac{P_\theta^\Pi(T_n)}{P_{\theta'}^\Pi(T_n)}
 =\prod_{t=1}^n\frac{K_{A_t}(V_t\mid\theta)}{K_{A_t}(V_t\mid\theta')}.
\]
具有零选择概率的路径对两者同时为零，不影响等式在实际支持上的使用。

\begin{theorem}[自适应相对熵账本]
\label{thm:kl-ledger}
令$N_a(n)$为实验$a$的调用数，令
$D_a(\theta\Vert\theta')=\KL(K_a(\cdot\mid\theta)\Vert K_a(\cdot\mid\theta'))$。则
\[
 \KL(P_\theta^\Pi\Vert P_{\theta'}^\Pi)
 =\sum_a\EE_\theta^\Pi N_a(n)D_a(\theta\Vert\theta').
\]
\end{theorem}
\begin{proof}
对似然比取对数再取$\theta$下的期望。条件于过去与当前实验，单项对数比的期望为$D_a$。把每一步按其所选实验求和即得。提前停止后的空操作在两个情境下相同，贡献零。
\end{proof}

这样的换测度方法是序贯实验与主动检验的基本工具\citep{chernoff1959,kaufmann2016}。这里的等式允许任意历史依赖的选取规则；重复、分叉、合并与多Agent产生的全记录都在比较范围内。

\section{正确回答要求多少相对熵}
定义二元相对熵
\[
 \kl(p,q)=p\ln\frac pq+(1-p)\ln\frac{1-p}{1-q}.
\]
若两个情境要求不同目标，而策略在每个情境的错误都不超过$\delta$，取事件$E=\{\widehat Y=g(\theta)\}$，则$P_\theta(E)\ge1-\delta$、$P_{\theta'}(E)\le\delta$。将完整记录压缩为事件是否发生，数据处理给出
\begin{equation}
 \sum_a\EE_\theta N_aD_a(\theta\Vert\theta')
 \ge\kl(1-\delta,\delta)
 =(1-2\delta)\ln\frac{1-\delta}{\delta}.
 \label{eq:pair-lower}
\end{equation}
所用二元单调性与数据处理在附录\ref{app:tools}证明。

\section{允许针对不同历史重新分配}
\begin{definition}[局部KL参数]
\[
 \kappa_{\mathrm K}(\theta)=
 \max_{w\in\Delta(\calA)}\min_{\theta'\in\Alt(\theta)}
       \sum_aw_aD_a(\theta\Vert\theta'),\qquad
 \kapK=\min_\theta\kappa_{\mathrm K}(\theta).
\]
\end{definition}
内层分配可以随所分析的真实情境变化。这给自适应组织比固定全局分配更大的自由度，并没有向实际控制者告知真实情境。

\begin{theorem}[全类别调用下界]
\label{thm:adaptive-lower}
若$\kapK>0$，任意最坏错误率不超过$\delta$的组织都满足
\[
 n\ge\frac{\kl(1-\delta,\delta)}{\kapK}.
\]
若$\kapK=0$，则存在目标冲突的不可辨别对。
\end{theorem}
\begin{proof}
固定$\theta$，令$q_a=\EE_\theta N_a/n$。其和不超过1，可将剩余质量补给任意实验而不减小区别量。式\eqref{eq:pair-lower}对所有冲突情境都成立，所以$n\kappa_{\mathrm K}(\theta)\ge\kl(1-\delta,\delta)$。对$\theta$取最严格条件即得。最后，有限多约束下严格正的全支持实验分配能区分所有非等同分布，故参数为零恰有冲突对的全部KL为零。
\end{proof}

\section{有了适应性是否一定更强}
下界的局部分配与上界的全局分配不同。某些情境中，先取得少量定位信息，再专门诊断当前备选集可能很有效；有些实验族则无需这种调整。主动序贯检验已系统研究此类选择\citep{naghshvar2013}。

两个不同参数出现在上下界中，表示尚待分析的结构差距。不能直接将$\kapK$当作已经可达的固定预算最坏情境误差指数。下一章在明确条件下连接二者；第\ref{ch:frontiers}章保留增长任务族中的适应性问题。

\begin{exercise}证明$\kl(p,q)$在$p\ge1-\delta$、$q\le\delta$、$\delta<1/2$时不小于$\kl(1-\delta,\delta)$。\end{exercise}
\begin{exercise}某策略在第一次观察为1时停止，否则继续。解释为何停止事件也携带信息，以及补空操作的记录为何没有丢掉它。\end{exercise}

\chapter{覆盖率的操作性含义}
\label{ch:operational}
\chapterguide{把独立定义的结构量同时连接到可实现构造和任意组织的下界，才得到关于组织能力的可操作结论。}

\section{同一个量何时控制两端}
假设存在$M\ge1$，对所有实验、片段和情境对，
\begin{equation}
 M^{-1}\le\frac{K_a(v\mid\theta)}{K_a(v\mid\theta')}\le M.
 \label{eq:ratio}
\end{equation}
有限严格正表总存在某个$M$；比较一族实例的数量级时，必须要求它们共享同一个上界，不能把随规模爆炸的$M$当作常数。

\begin{lemma}[KL与Hellinger区别的比较]
\label{lem:comparison}
满足式\eqref{eq:ratio}时，
\[
 D_a(\theta\Vert\theta')\le C_Mh_a(\theta,\theta'),
 \qquad C_M=2(\sqrt M+1)^2.
\]
\end{lemma}
\begin{proof}
由$\ln x\le x-1$，$\KL(P\Vert Q)\le\sum_v(P-Q)^2/Q$。而
\[
 \frac{(P-Q)^2}{Q}=(\sqrt P-\sqrt Q)^2(\sqrt{P/Q}+1)^2
 \le(\sqrt M+1)^2(\sqrt P-\sqrt Q)^2.
\]
求和，并使用平方根距离等于$2h$。
\end{proof}

\begin{proposition}[局部分配到全局分配]
\label{prop:local-global}
情境数为$d$时，$\kapK\le dC_M\kap$。
\end{proposition}
\begin{proof}
为每个$\theta$取达到$\kappa_{\mathrm K}(\theta)$的分配$w^\theta$。其对每个目标冲突情境的平均Hellinger区别至少为$\kapK/C_M$。将这$d$份分配等权混合。任意冲突对至少从其中与其一个端点对应的分配获得$\kapK/(dC_M)$，故全局最优覆盖率不小于此值。
\end{proof}

\begin{theorem}[有限诊断组织的操作性界]
\label{thm:main}
对已知有限实验表、至少两个目标、有界似然比和单位调用成本，若$\kap>0$，则
\begin{equation}
 \frac{\kl(1-\delta,\delta)}{dC_M\kap}
 \le n^*(\delta)
 \le\left\lceil\frac{\ln(m/\delta)}{\kap}\right\rceil.
 \label{eq:main}
\end{equation}
若$\kap=0$，则任意$\delta<1/2$都不能得到最坏情境保证。固定$d,M$并令$0<\delta\le1/4$时，可辨别实例满足
\[
 n^*(\delta)=\Theta_{d,M}\!\left(\kap^{-1}\ln(1/\delta)\right).
\]
\end{theorem}
\begin{proof}
下界连接定理\ref{thm:adaptive-lower}与命题\ref{prop:local-global}，上界来自定理\ref{thm:upper}。零覆盖率情形来自不可辨别性。对$\delta\le1/4$，$\kl(1-\delta,\delta)$与$\ln(1/\delta)$之比有统一正下界；又有$m\le d-1$。上界的取整可吸收到常数中，因$\kap\le1$且$\ln(1/\delta)\ge\ln4$。
\end{proof}

\section{这个结论实际解释什么}
覆盖率从各实验产生的分布计算，它并未把最优求解成本写进定义。主定理说明，在固定有限情境与似然比条件下，它的倒数给出降低错误所需调用数的尺度。组织的作用是将基本操作中分散的区别转化为统一的任务可靠性。

这里的比较方包含任意自适应调度、任意中间表示和任意计算能力的控制者。所证明的下界不依赖把对手限制为整条重试；所构造的上界也不要求检查器在每次历史下预先保有可靠余量。

代价同样清楚：主定理只控制固定有限任务模型的统计调用复杂度。$d$增长时，下界损失一个显式$d$因子；上界可能为辨认无关的情境细节多付出$\ln m$。因此覆盖率尚未独自完成大规模任务的最优复杂度刻画。

\section{不能从覆盖率中消去任务规模}
令$\theta\in\mathbb F_2^r$，目标是恢复整个$\theta$。只有一种实验：每次独立产生均匀随机向量$U\in\mathbb F_2^r$，返回片段$(U,U^\top\theta\oplus E)$，其中$E\sim\Bern(\eta)$，固定$0<\eta<1/2$。向量$U$是实验输出的一部分，控制者不能免费选择它。

\begin{theorem}[相同覆盖率下的规模增长]
\label{thm:dimension}
该实例族对所有$r\ge1$都具有
\[
 \kap=\tfrac12\bigl(1-2\sqrt{\eta(1-\eta)}\bigr),
 \qquad M=(1-\eta)/\eta.
\]
二者均与$r$无关。但任何最坏错误不超过$\delta$的组织仍需
\[
 n\ge\frac{(1-\delta)r-h_2(\delta)}{1-h_2(\eta)}.
\]
固定$\eta$、$0<\delta\le1/4$时，调用复杂度为
$n^*(\delta)=\Theta_\eta(r+\ln(1/\delta))$。
\end{theorem}
\begin{proof}
任意非零情境差$z$满足$U^\top z=1$的概率为$1/2$。在另一半输出向量上，两种bit分布相同；在这一半上它们互换。因此平均区别为$h_\eta/2$，概率比始终位于$[M^{-1},M]$。

对均匀情境先验，目标熵为$r$ bit。每次的新$U$独立于情境与历史；给定$U$后，带噪bit对情境的条件互信息至多$1-h_2(\eta)$，因为其熵至多1，而条件于情境的噪声熵为$h_2(\eta)$。信息链式法则和Fano界给出所写下界。

当$r\ge2$，$\delta\le1/4$时分子至少$r/4$。另一方面，任意情境对的单次KL均为$d_\eta/2$，定理\ref{thm:adaptive-lower}再给出关于$\ln(1/\delta)$的下界；这也覆盖$r=1$。两个下界的最大值至少为其和的一半。上界由$m=2^r-1$和定理\ref{thm:upper}得到。
\end{proof}

这个例子同时说明覆盖强度与需要恢复的区别数量各有作用。没有发生模型失配或似然比爆炸，单独知道覆盖率仍无法控制增长任务族的预算。实验输出长度为$r+1$ bit，其传输与读取成本还应另外计入。

\section{完整执行账本}
一个直接实现先计算概率表中的区别量与覆盖分配，再运行$n$次调用，在线累计$d$个候选的对数似然。若单次调用成本为$c_F$、一个数值更新成本为$c_{\mathrm{num}}$，则可写成
\[
 C_{\mathrm{total}}=C_{\mathrm{table}}+C_{\mathrm{design}}
       +nc_F+O(nd)c_{\mathrm{num}}+C_{\mathrm{communication}}.
\]
数值项计入实际精度下的加法与查表费用。概率表描述长度、求解规划以及控制精度都进入前置费用；表规模指数大时，这种实现可能不经济。

\section{一个充分结果怎样变成研究工具}
可以用主定理检查组织建议。若建议只复制完全相同的实验分布，它不会创造新的可辨别对。若建议引入能打破旧混淆类的操作，可先检查覆盖率是否从零变正。若各实验原本都可用，建议只是调整频率，则可比较分配的最差覆盖值。后两种改变分别涉及允许能力的改变与固定能力内的组织改进，报告时需要区分。

\begin{exercise}说明为什么$M$在每个有限实例中存在，不足以证明一族规模增长实例具有统一常数的复杂度界。\end{exercise}
\begin{exercise}若已找到覆盖率$u>0$的分配，但不知道最优值，仍能作出哪些可验证的可靠性与成本承诺？\end{exercise}

\part{共同错误与通信}
\chapter{共同偏差、校准与可构造表示}
\label{ch:calibration}
\chapterguide{通过有限域上的观察关系，可以判断共同错误何时能够消除，并实际构造所需的组合规则。}

\section{把持续失效因素写入情境}
令$\theta\in\mathbb F_2^r$，所有情境都可能出现。基本实验由向量$a\in\calA\subseteq\mathbb F_2^r$标记，输出
\[
 V=a^\top\theta\oplus E,\qquad E\sim\Bern(\eta),\quad0<\eta<1/2.
\]
各次$E$独立，未知$\theta$始终相同。目标为$g(\theta)=l^\top\theta$，其中$l\ne0$。一个坐标可以表示真实事实，也可以表示影响多种规则的共同偏差；每行$a$说明一个操作把哪些因素混在一起。

\begin{theorem}[线性目标的可辨别性]
\label{thm:span}
目标允许任意小的最坏错误率，当且仅当
\[
 l\in\operatorname{span}_{\mathbb F_2}(\calA).
\]
\end{theorem}
\begin{proof}
两个情境之差为$z$时，实验分布相同恰当$a^\top z=0$，因为$\eta\ne1/2$。因此不可辨别的差组成全部实验行的零空间。目标在类上为常数恰当$l^\top z=0$对该零空间中的每个$z$成立；有限维非退化内积下，这等价于$l$属于行空间。也可直接通过消元构造$l$的线性组合，或找到零空间中满足$l^\top z=1$的反证向量。最后应用推论\ref{cor:identifiability}。
\end{proof}

\section{区别强度与结构覆盖分离}
若$a^\top z=0$，两分布相同；否则它们分别为参数$\eta$与$1-\eta$的Bernoulli分布。因此
\[
 h_a(\theta,\theta')=h_\eta\,\mathbf1\{a^\top(\theta-\theta')=1\},
 \qquad h_\eta=1-2\sqrt{\eta(1-\eta)}.
\]
定义
\[
 \mu(\calA,l)=\max_w\min_{z:\,l^\top z=1}
               \sum_aw_a\mathbf1\{a^\top z=1\}.
\]
于是$\kap=h_\eta\mu$。局部KL分布在所有真实情境下具有同样的冲突差集合，又有
\[
 \kapK=d_\eta\mu,\qquad
 d_\eta=(1-2\eta)\ln\frac{1-\eta}{\eta}.
\]
该模型的下界可以直接使用$d_\eta\mu$，无需付出主定理的一般$dC_M$比较因子。基本实验的噪声强度与它们覆盖哪些错误方向，在这里成为两个独立可计算因素。

\section{校准一个共同偏差}
取$\theta=(Y,B)$、$l=(1,0)$。原实验$a_1=(1,1)$观察$Y\oplus B$，校准实验$a_2=(0,1)$观察$B$。只有$a_1$时目标不在行空间内，均匀先验的误差下限为$1/2$。两种实验都允许时，$l=a_1+a_2$。

目标冲突差只有$(1,0)$与$(1,1)$。前者仅由$a_1$区分，后者仅由$a_2$区分，故
\[
 \mu=\max_{w\in[0,1]}\min(w,1-w)=1/2,
 \qquad \kap=h_\eta/2.
\]
两个实验都已经属于基本能力时，开始使用校准实验属于组织改变；若原系统根本没有该操作，增添它则改变了允许的基本能力集合。两者都可能消除盲点，但归因不同。

\begin{figure}[ht]
\centering
\begin{tikzpicture}[x=6cm,y=4cm]
\draw[edge] (0,0)--(1.1,0) node[right] {$w$};
\draw[edge] (0,0)--(0,0.65) node[above] {覆盖率/$h_\eta$};
\draw[ink,thick] (0,0)--(0.5,0.5)--(1,0);
\draw[dashed,accent] (0.5,0)--(0.5,0.5)--(0,0.5);
\node[below] at (0.5,0) {$1/2$};
\node[left] at (0,0.5) {$1/2$};
\node[below] at (1,0) {$1$};
\node[above right,font=\small] at (0.5,0.5) {均衡分配};
\end{tikzpicture}
\caption{校准模型的最差诊断覆盖率。$w$为使用原观察的比例。}
\end{figure}

\begin{table}[ht]
\centering
\begin{tabular}{rrrr}
\toprule
噪声 $\eta$ & 覆盖率 & 必要调用下界 & 充分调用预算\\
\midrule
0.1 & 0.2000 & 4 & 19\\
0.2 & 0.1000 & 7 & 37\\
0.3 & 0.0417 & 16 & 89\\
0.4 & 0.0101 & 66 & 366\\
\bottomrule
\end{tabular}

\caption{校准模型在$\delta=0.05$下的解析预算界。下界约束所有策略，上界由公开随机分配实现；二者不表示有限参数下已经达到精确最优。}
\label{tab:calibration}
\end{table}

\section{表示怎样实际构造}
对实验行做二元高斯消元，可以在关于显式矩阵尺寸的多项式时间内判断$l$是否可表达。若可表达，就得到
$l=a_1+\cdots+a_k$，去掉重复行后可令$k\le r$。目标等于这$k$个真实观察bit的异或。

\begin{proposition}[已构造线性表示的可靠执行]
\label{prop:parity-execute}
取奇数$s\ge\ln(k/\delta)/(2(1/2-\eta)^2)$，每个所选实验调用$s$次并取多数，最后异或这些多数结果。调用$ks$次，最坏错误率不超过$\delta$。
\end{proposition}
\begin{proof}
每个多数错误概率至多$e^{-2s(1/2-\eta)^2}\le\delta/k$，由独立有界和的集中界。全部$k$项正确时异或必正确，故并合界给出结论。消元、计数与异或不需要读取未知情境。
\end{proof}

这个构造解决了本模型中表示如何取得的问题。它可能比枚举$2^r$个情境计算似然经济得多。若实验只有标准基且$l$包含$k$个坐标，均匀分配给这$k$行给出$\mu=1/k$；反方分别选择这$k$个单坐标差，证明无法更大。一般行集合中的分配由有限线性规划确定。

\begin{exercise}对$\calA=\{(1,1,0),(0,1,1)\}$，判断目标$(1,0,1)$与$(1,1,1)$各自是否可辨别，并给出组合或反证。\end{exercise}
\begin{exercise}证明标准基实验、含$k$个非零坐标的目标具有$\mu=1/k$。\end{exercise}

\chapter{通信：为另一个过程生成什么}
\label{ch:communication}
\chapterguide{消息的作用由接收者仍需区分的情境决定。生成、公开与采用分别发生，三者共同决定证据怎样进入组织。}

\section{局部历史中的通信选择}
Agent A已经检查了一部分证据，Agent B正决定是否继续某项方案。A可以转发已有观察，也可以专门生成一份与B的备选解释有关的似然摘要；后者可能在准备前尚未存在。

令$h_t^A$为A实际知道的历史。通信规则$\Gamma_A(h_t^A)$选择接收者、消息生成操作，以及立即公开、继续准备或等待。生成新片段仍是一条$\II_e\xrightarrow{F_a}v$；使片段对B可见是另外的事件。B的模式决定是否采用它，收到猜想不等于把猜想当作有效前提。

控制者的全局实验表不能代替A的局部信息。A可以依据公开计划或B发来的请求了解其需求。无论规则由谁生成，它在运行时都只能使用已经实际可见的材料。

\section{通信怎样改变诊断实验}
设某次实验的原片段为$v$，消息按与真实情境无关的核$Q(m\mid v)$生成。接收者看见的实验为
\[
 \widetilde K_a(m\mid\theta)=\sum_vQ(m\mid v)K_a(v\mid\theta).
\]
这是统计实验的后处理\citep{blackwell1953}。消息记录必须包含接收者能观察到的全部选择信号，适用时也包括发送标记与时刻。

\begin{theorem}[通信压缩的覆盖单调性]
\label{thm:garbling}
对每个实验和情境对，$h(\widetilde K_a^\theta,\widetilde K_a^{\theta'})\le h(K_a^\theta,K_a^{\theta'})$。因此，在相同实验分配集合下，$\widetilde\kap\le\kap$；不可辨别类只能合并或保持。
\end{theorem}
\begin{proof}
对每个消息$m$，Cauchy--Schwarz给出
\[
 \sqrt{\sum_vQ(m\mid v)P(v)\sum_vQ(m\mid v)P'(v)}
 \ge\sum_vQ(m\mid v)\sqrt{P(v)P'(v)}.
\]
对$m$求和，亲和度不会减小，区别量不会增大。逐项取加权和、最小值与最大值得到覆盖结论；原来相同的分布经过同样后处理仍相同。
\end{proof}

这不意味着压缩无用：传递与解码成本可能显著下降。定理提供必须一起核对的另一面——节约是否毁掉了任务依赖的区别。

\section{什么压缩完整保留了证据}
先考虑严格正概率表和确定性编码$M=\phi(V)$。某个编码单元是产生同一消息的片段集合。

\begin{proposition}[亲和度无损压缩]
\label{prop:sufficiency}
对一对情境，压缩前后亲和度相等，当且仅当每个编码单元中$K_a(v\mid\theta)/K_a(v\mid\theta')$为常数。对全部情境对无损，等价于每个单元中的似然向量$\bigl(K_a(v\mid\vartheta)\bigr)_{\vartheta\in\Theta}$互相成正比例。
\end{proposition}
\begin{proof}
确定性编码下，每个消息对应一个求和单元。上一定理的每个Cauchy--Schwarz差非负，和为零恰当各单元取等。严格正条件下，取等恰当平方根概率向量成比例，等价于所写似然比条件。对全部情境对应用该条件得到向量表述。
\end{proof}

无损压缩因此可以从接收者的备选模型出发设计。表达形式可以变化，只要这些似然区别被保留，统计证据就没有减少。保留某个任务当前决策所需的信息，有时还允许更强压缩；需要另行指定损失、先验与后续问题。

\section{为接收者生成一份充分记录}
允许实验安排根据已有观察自适应改变，但选择规则只能使用记录和情境独立随机性。$J$个成员在汇总点提交一份完整前缀的计数：每项携带情境信息的观察都恰好计入一个成员，不能遗漏用于决定下一实验的私有观察。第$j$个成员统计
\[
 C_{j,a,v}=\#\{t:\text{该调用属于成员 }j,\ A_t=a,V_t=v\}.
\]
汇总者计算$\ln L_n(\vartheta)=\sum_{j,a,v}C_{j,a,v}\ln K_a(v\mid\vartheta)$，便与接收全部原始记录得到相同的情境似然比和最大似然输出。

\begin{proposition}[计数消息的精确组织实现]
\label{prop:counts}
上述完整计数保留任意合法自适应组织的全部情境似然比。若总调用数不超过$n$，其一个固定编码使用至多
\[
 J|\calA||\VV|\lceil\log_2(n+1)\rceil
\]
bit，另加约定的传输封装成本。
\end{proposition}
\begin{proof}
任意完整记录的概率可写成$q_\Pi(T)\prod_{j,a,v}K_a(v\mid\vartheta)^{C_{j,a,v}}$，其中$q_\Pi(T)$是策略在该路径上的选择概率乘积，与$\vartheta$无关。将同一计数$C$的记录求和，得到
\[
 \PP_\vartheta(C)=\left(\sum_{T:\,C(T)=C}q_\Pi(T)\right)
       \prod_{j,a,v}K_a(v\mid\vartheta)^{C_{j,a,v}}.
\]
括号因子对全部情境相同，在似然比中消去。每个计数为$0$到$n$的整数，固定编码各用$\lceil\log_2(n+1)\rceil$ bit足够。
\end{proof}

准备这样的片段使接收者不必重放整个观察块。自适应选择本身不会破坏完整计数的充分性。困难发生在只统计被筛选公开的观察、漏掉决定筛选的私有证据时；此时应把缺少的证据计入，或推导包含选择与沉默的消息似然。充分性也不自动保证任意调度具有良好覆盖率，可靠性上界仍需另外满足相应分配条件。

\section{何时证据已经足够}
考虑接收者拥有完整证据前缀的情形。概率核严格正，选实验的规则可以自适应。令$\ell_t(\vartheta)=\ln L_t(\vartheta)$，阈值$b=\ln(m/\delta)$。首次存在候选$\vartheta$满足
\begin{equation}
 \ell_t(\vartheta)-\ell_t(\vartheta')\ge b
 \quad\text{对全部 }\vartheta'\in\Alt(\vartheta)
 \label{eq:stop}
\end{equation}
时输出$g(\vartheta)$并停止。规则可以在固定批次消息到达后才检查。

\begin{theorem}[随时有效的停止规则]
\label{thm:stop}
式\eqref{eq:stop}在任意实验选择策略下，作出错误停止的概率不超过$\delta$。若持续使用覆盖率$u>0$的独立分配，并在每个前缀检查，则停止时刻$\tau$几乎必然有限，且
\[
 \PP_\theta(\tau>n)\le\min\{1,me^{b/2}(1-u)^n\}.
\]
\end{theorem}
\begin{proof}
固定真实$\theta$。对每个冲突情境$\theta'$，$L_t(\theta')/L_t(\theta)$是初值1的非负鞅。若错误选中$\theta'$，该比值必曾达到$e^b$。非负鞅最大不等式给出概率至多$e^{-b}$；对至多$m$个冲突情境求和为$\delta$。

对独立分配，若在$n$时所有真实情境对冲突情境的对数似然差都达到$b$，必已停止。任意冲突情境违反该条件的概率，由平方根似然比的Markov不等式，至多$e^{b/2}(1-u)^n$。求和得到尾界，故停止几乎必然发生。
\end{proof}

随时有效的界通过鞅而非反复使用固定样本界成立\citep{howard2020}。仅仅在很多时刻重复检查某个固定置信阈值，通常不能自动保持总错误预算。

\section{沉默与选择性公开}
若A知道一个公平bit，双方约定可靠时隙，A仅在bit为1时发送标记，则B从“标记或沉默”完整恢复一bit。平均只有半次显式发送，但可见记录仍有一bit熵。发送次数、消息长度和完整记录信息量必须分别计算。

更一般地，以$S$表示发送与否，$V$在未发送时取固定空值，则
\[
 I_2(Y;S,V\mid Z)=I_2(Y;S\mid Z)
  +\PP(S=1)I_2(Y;V\mid Z,S=1).
\]
$Z$是接收者已有材料。等式由条件链式法则得出。忽略第一项，就会遗漏选择性公开本身的证据。若B不知道A是否已经完成推理，沉默的含义也随之改变。

\section{消息信息与任务损失}
接收者已有$Z$，需要预测目标$Y$，消息为$M$。最小平均对数损失由$H_2(Y\mid Z)$降为$H_2(Y\mid M,Z)$，差值为$I_2(Y;M\mid Z)$；条件交叉熵分解直接证明这一点。

若目标字母表大小为$q\ge2$，重构错误率至多$\delta\le1-1/q$，则条件Fano界给出
\[
 I_2(Y;M\mid Z)\ge H_2(Y\mid Z)-h_2(\delta)-\delta\log_2(q-1).
\]
完整消息用$L$个固定二元槽位表达时，左边至多$L$。信息与编码的基本联系源于通信理论\citep{shannon1948}；这里的诊断覆盖率进一步分析消息是否保留所有目标冲突对的区别。

\section{局部请求怎样改变内容}
若B只需知道某个候选是否被当前证据排除，A可以生成相应似然差；若B还要在许多候选之间选择，则完整似然向量或计数可能更适合。接收者的请求改变待保留的区别，从而改变应该执行的消息生成操作。

公开时机可以以式\eqref{eq:stop}所需的证据为依据，也可以按固定预算分批。更一般的通信价格、延迟与私有状态联合最优控制超出这个停止定理；但任何方案都必须把可见的沉默、筛选和时间信号放进同一个观察模型。

\begin{exercise}为什么将实验输出全部压缩成一个固定符号，会使任何诊断覆盖率变为零？\end{exercise}
\begin{exercise}给出一个选择性发送协议，说明只保留已发送内容的似然而忽略发送规则，会算错接收者后验。\end{exercise}

\chapter{表示、计算与不准确的模型}
\label{ch:representation}
\chapterguide{保存区别、低成本读出区别，以及正确描述区别，是三个需要分别证明的条件。}

\section{组织不能把未知答案藏进新操作}
假设将$n$步协议的最终输出重新命名为一次宏操作。它形成一个新概率核，但执行这个核仍需原来的$n$次调用。将宏操作当作单位成本，会人为改变能力比较。

对任意两情境，宏输出的KL不大于完整记录的KL，而完整记录KL等于第\ref{ch:lower}章的调用账本。因此后处理可以降低计算或通信负担，不能绕过原有证据下界。若某种“检查”真的提供了此前没有的观察分布，它应作为新基本能力计费。

\section{同样的信息可以需要不同读出成本}
考虑已经观察到的$k\ge2$个独立公平bit，目标是它们的异或。一个只能读取一个原坐标的读出器，错误率为$1/2$；先计算异或并把结果放入一个新片段，同一类单坐标读出器便能零错误读取目标。

原数据与新片段都含有一bit目标信息。变化发生在允许的读出计算，而异或计算已经由准备片段的过程付出$O(k)$的费用。可利用信息理论明确将预测器的能力纳入度量\citep{xu2020}。这里的诊断定理首先给出统计限制，实现时再列出取得相应解码器所需成本。

第\ref{ch:calibration}章的消元构造具有不同作用：它在显式实验矩阵上找到目标的表达，再用有噪观察执行表达。由此避免直接枚举$2^r$个情境。表示的收益通过可执行算法体现，并非仅给原问题换一个符号。

\section{概率表估计错误会消耗多少余量}
实际解码使用估计表$\widehat K$，真实表仍为$K$。设所有概率严格正，并有统一对数误差
\[
 |\ln\widehat K_a(v\mid\vartheta)-\ln K_a(v\mid\vartheta)|\le\varepsilon.
\]
分配$w$在真实表上具有最差覆盖率$u$。估计最大似然输出按$\widehat L$选择情境。

\begin{theorem}[模型误差的可靠性代价]
\label{thm:robust}
上述组织满足
\[
 \max_\theta\PP_\theta(g(\widehat\theta)\ne g(\theta))
 \le\min\{1,me^{n\varepsilon}(1-u)^n\}
 \le\min\{1,me^{-n(u-\varepsilon)}\}.
\]
\end{theorem}
\begin{proof}
每个候选的累计对数似然误差至多$n\varepsilon$。若错误候选在估计表上胜出，则真实对数似然差$\ln L_n(\theta')-\ln L_n(\theta)$至少为$-2n\varepsilon$。该事件的概率至多$e^{n\varepsilon}$乘两种真实记录分布的亲和度。再对至多$m$个候选求和。
\end{proof}

若$u>\varepsilon$，这个保证仍给出指数压低误差的构造。若$u\le\varepsilon$，所写上界不再保证收敛；这不等价于已经证明所有算法失败。统一对数误差需要控制小概率项，有限样本平均拟合误差无法自动替代它。

同一证明也给通信量化界：若接收者累计的每个对数似然分数误差至多$E$，最大似然风险至多$me^E(1-u)^n$。这允许根据可用诊断余量选择摘要精度。仅仅传递浮点分数而不规定误差，会使所谓“充分消息”失去保证。

\section{模型失配的更深层问题}
上述定理要求真实情境已在候选集合中。若全部候选共享错误解释，增加观察会可靠地选择一个失配模型，不能得到关于真实目标的保证。

对一组可能概率表，可先检查它们是否包含两个目标不同但全部观察分布相同的实例。一旦存在这样的对，定理\ref{thm:blind}直接给出稳健推断的障碍。不存在这类对也未必意味着容易学习：还要解决实验设计、有限样本辨认与模型表示成本。

\begin{exercise}每个成员提供的各候选对数似然分数误差至多$e_j$。给出汇总后风险上界，并说明误差预算如何相加。\end{exercise}
\begin{exercise}如果原始片段的内容只经过可逆重编码，为什么信息区别不变，而实际解码时间仍可能改变？\end{exercise}

\part{解释与研究边界}
\chapter{认知现象的结构解释}
\label{ch:interpretation}
\chapterguide{一个解释需要指出什么结构发生了变化，以及这种变化怎样影响可验证的任务表现。概念名称本身不承担证明。}

\section{无知、混淆与过度确信}
给定任务与实验族，无知可以表现为当前证据尚不足；结构性混淆则指两个目标冲突情境属于同一不可辨别类。前者可能通过继续当前实验改善，后者需要改变可取得或可使用的区别。

校准模型给出这一区别的可观察后果：只重复带偏差观察，会提高对$Y\oplus B$的确定性，但关于$Y$的Bayes错误仍为$1/2$。增加校准操作以后，目标可辨别，适当分配能够继续压低错误。因此“更加确信”需要绑定正在确信的变量，不能仅根据推理长度判断。

\section{好奇、专注与无效重复}
在覆盖模型中，选择实验可以针对尚未充分区分的冲突对。已经取得很多证据的方向继续占用大部分预算，会让其他方向成为最坏情境的瓶颈。切换操作具有明确的诊断理由。

这可以解释某些探索与专注行为：探索寻找未覆盖的错误方向，专注积累对当前关键方向的区别。它们的合理性依赖任务、已有证据与允许操作。覆盖率是固定全局分配的尺度，自适应阶段切换的最佳规则还需要更精细的分析。

无效重复可以用两种机制说明。第一种是结构性零增益：操作对当前关键冲突分布完全相同。第二种是资源性低增益：区别虽正，但相对于调用或读出成本很弱。将“无聊”解释成某种停止倾向时，应说明采用哪一种机制，避免仅把短期没有进展重新命名。

\section{理解与可迁移的组合}
在二元线性模型中，一份可验证的表示$l=a_1+\cdots+a_k$允许系统在全部未知情境上组合局部观察。它支持跨实例的同一操作方式。这样的迁移保证比记住某个实例的答案强，并可以通过不同情境下的风险检查。

但如果任务目标变为另一个$l'$，原表示未必仍然可用。理解的范围应由所覆盖的目标集合描述。例如，实验行空间给出全部可辨别线性目标；构造短表示的算法与它的成本，再给出实际可用的范围。

\section{发现、创新与新的可用区别}
一项改变至少可能发生在三个层次。发现新的实验，可能细分不可辨别类；在既有操作中找到更好的组织，可能改善诊断分配；找到低成本表示，可能保留信息同时减少取得答案的计算。三者各有不同的资源来源。

若要评价一个改变相对于旧系统的价值，可以固定任务族和预算，检验错误下限是否下降、达到指定可靠性所需调用是否减少、或计算与通信费用是否降低。新片段的表面差异不直接进入这些比较。

这里得到的是几种可分析的创新机制。一个具体变化是否属于其中之一，需要提供操作、概率或程序上的证据；“创新”这个名称无法代替这些结构。

\section{社会组织与成员差异}
两个成员使用不同措辞，但全部实验分布相同，它们的存在不会打破不可辨别类。成员使用互补实验，或一个成员能够把另一个成员难以解码的记录变成低成本充分消息，则可能改善整体系统。

对于同一完整信息与计算条件，一个集中控制者可以模拟有限团队。多成员优势需要说明额外资源：分散的观察权限、并行计算、局部专业读出，或者更易实现的模式。通信章的充分计数实现展示了一个明确情形：分散采样与集中似然判断可以得到相同的统计保证，同时只传递所需记录。

\begin{exercise}提出一个可反驳“更多独立成员必定消除共同错误”的有限实验，并说明独立性的条件是什么。\end{exercise}
\begin{exercise}对一个新的表示，分别设计信息保持、计算成本和跨实例可靠性三项检查。\end{exercise}

\chapter{从有限定理走向更强理论}
\label{ch:frontiers}
\chapterguide{进一步研究应当改进核心模型的可达范围与解释力，并保留明确的反例和比较边界。}

\section{当前已经连接起来的结果}
不可辨别类决定永久误差，覆盖率决定一个可执行组织的可靠性，自适应KL账本约束全部策略。在固定$d,M$下，三者形成操作性主定理。校准模型给出表示存在性与构造算法；通信后处理、充分计数和量化界说明组织改变何时保存这些保证。

这条链采用经典统计实验、主动检验、信息不等式和有限线性代数。其价值应根据模型是否清楚、连接是否正确、解释是否可复用来判断。哪些结果能形成独立原创贡献，还需要逐项与最接近的研究比较；此处没有把重新命名的经典量列作发现。

\section{问题一：规模增长时，覆盖率还缺少什么}
主定理含$d$因子，上界含$\ln m$。当情境数随问题规模指数增长时，这些因素可能压倒结论。需要研究目标冲突图、概率表的对称性和可构造表示，能否替代原始情境数。

二元模型提供一个具体测试：标准基实验与$k$位异或目标具有$\mu=1/k$，局部多数构造需要$O(k\ln(k/\delta)/(1/2-\eta)^2)$次调用，而通用最大似然并合界可能更松。能否对完整自适应类别证明更紧下界，并找出决定额外对数因子的结构？这里可以检验覆盖率是否需要新的互补参数。

\section{问题二：什么时候适应性具有不可替代收益}
全局覆盖分配同时保护全部情境，局部KL规划允许对真实情境分别设计。要研究两者之间的实际差距，需要把学习当前所在区域的成本加入策略，并比较最坏风险、平均风险、固定预算和期望停止时间等不同标准。

一个有内容的分离应同时给出：所有非自适应组织的下界、可执行自适应组织的上界，以及不依赖新增观察权限的比较。不能仅凭两个形式不同的规划值就宣布可达的适应性优势。

\section{问题三：共同错误在过程中变化}
静态$\theta$表达持续的共同因素。若一次交谈改变其他成员的错误状态，或一次操作破坏未来观察条件，输出核需要依赖隐藏状态与历史。组织此时同时选择取证与改变情境。

可以先限定有限状态、已知转移和局部可见性，研究是否存在任何能把两个目标冲突状态区分开的控制序列。静态模型中的一次实验混淆类将变成受控可辨别性问题。新的上下界应把等待、干预和错误传播计入状态转移。

\section{问题四：观察模型怎样取得}
已知概率表把研究集中在组织。未知表需要从任务与诊断样本中学习，并保证对小概率项、共享偏差和未见情境的处理。第\ref{thm:robust}条定理只说明统一对数误差如何消耗余量，尚未提供取得该误差保证的通用算法。

一个可推进的任务是：固定具有有限描述的概率表族，联合控制模型识别成本与之后的诊断成本，寻找先学习再组织与边学边组织的分界。允许任意未知表时，未观察区域上的不同解释可能永远不可辨别。

\section{问题五：通信怎样内生改变组织}
完整计数可以保存自适应组织的全部证据，选择性通信则未必交付完整计数。私有信息驱动的消息既改变接收者后验，又影响其他成员下一步实验。联合优化需要同时描述消息生成、可见沉默与局部策略，并对全部相关成本给出界。

可先研究两个成员、有限观察与有代价的可靠时隙：比较固定轮询、接收者请求和事件触发的最优任务风险。若某种协议取得优势，需要证明优势来自更有效地利用已有局部信息，而未在发送时刻或控制信号中隐藏免费信道。

\section{研究成果的检验单位}
一项结果应能独立交付模型、定理、证明、反例与最接近文献的差异。对一般性结论，数值核验只负责发现有限实例中的错误；对实际Agent，则需另外估计实验模型并检验失配。

可以优先解决上述某一个窄问题，再看新结构是否能回过头改进主定理、校准算法和通信分析。由同一结果推动多个部分，才能使理论继续统一，而非依靠增加概念扩大篇幅。

\begin{exercise}选择一个增长的有限实例族，指出主定理中哪个常数随规模变化，并说明它是否仍给出有用界。\end{exercise}
\begin{exercise}为一个候选研究结果列出最小可审查材料：模型、上下界、反例、可复算实例和最近已有结果。\end{exercise}

\appendix
\chapter{概率、信息与对偶工具}
\label{app:tools}

\section{事件数据处理}
对于分布$P,Q$和事件$E$，将和式分别限制到$E$与其补集，条件相对熵的非负性给出
\[
 \KL(P\Vert Q)\ge\kl(P(E),Q(E)).
\]
更一般的确定性分组与随机后处理，可以通过加入共同的随机核再条件分解得到。对二元相对熵求导，
\[
 \partial_p\kl(p,q)=\ln\frac{p(1-q)}{q(1-p)},\qquad
 \partial_q\kl(p,q)=\frac{q-p}{q(1-q)}.
\]
当$p>q$时，前者为正、后者为负。因此$p\ge1-\delta$、$q\le\delta$时最小值在边界取得。

\section{独立有界和的集中}
令$X\in[0,1]$，$f(\lambda)=\ln\EE e^{\lambda(X-\EE X)}$。有$f(0)=f'(0)=0$，二阶导为指数倾斜分布中的方差，至多$1/4$。两次积分得$f(\lambda)\le\lambda^2/8$。

对独立的$X_1,\ldots,X_s$，指数矩相乘。Markov不等式与参数最优化给出
\[
 \PP\!\left(\frac1s\sum_i(X_i-\EE X_i)\le-\gamma\right)
 \le e^{-2s\gamma^2}.
\]
这就是本书多数纠错使用的有界集中工具\citep{hoeffding1963}。共同偏差不会满足以真实目标为中心的重复正确性假设；必须先由观察模型确定多数在估计哪个变量。

\section{非负鞅的随时界}
设$M_t\ge0$为初值1的鞅。令$\tau$为首次$M_t\ge c$的时刻。对有限$n$，有$\EE M_{\tau\wedge n}=1$，而在$\tau\le n$时该变量至少$c$。因此$\PP(\tau\le n)\le1/c$。令$n$趋于无穷，
\[
 \PP(\sup_tM_t\ge c)\le1/c.
\]
对情境似然比，给定过去和所选实验，其下一乘因子的条件期望为$\sum_vK_a(v\mid\theta')=1$，故确为鞅。这证明通信停止规则所需的全部条件。非负超鞅的同类论证见\citet{howard2020}。

\section{条件Fano界}
令$Y$有$q$种可能值，估计由记录$W$生成，错误指示为$E$。给定$W$，当$E=0$时目标确定；当$E=1$时最多有$q-1$个可能值。因此
\[
 H_2(Y\mid W)\le h_2(p_e)+p_e\log_2(q-1).
\]
右侧在$0\le p_e\le1-1/q$上递增，故可用该区间内的错误上界$\delta$替换。令$W=(M,Z)$并减去条件熵，就得到正文的信息下界。独立解码随机性可以并入$W$，不会额外提供目标信息。

\section{有限覆盖规划的对偶}
把每个实验写成一个向量，其坐标是各冲突对的区别量；令$C$为这些向量的凸包。覆盖率是$\max_{x\in C}\min_jx_j$。对任意坐标概率$\lambda$，有$\min_jx_j\le\lambda^\top x$，这先给出原问题不大于对偶。

取$t>\kap$。紧凸集$C$与上正交区域$U_t=\{x:x_j\ge t\}$不相交。有限维凸集分离定理给出一个严格分离方向$\lambda$。由于$U_t$沿每个正坐标无界，分离方向必须非负且不全零。归一化$\sum_j\lambda_j=1$后，
\[
 \max_{x\in C}\lambda^\top x<t.
\]
让$t\downarrow\kap$并用概率单纯形的紧性，得到对偶不大于原问题。于是二者相等。该论证只使用有限维凸分离，没有给数值求解器的精度附加假设。

\chapter{练习解答}

\section*{第\ref{ch:model}章}
1．一个数学上存在的随机变量未必属于可观察材料。让控制者读取真实种子或包含情境的变量，会改变可用信息与策略类别，从而取消原有下界的条件。

2．两个历史最后都得到“候选A仍可行”。第一个历史已经排除B，第二个尚未检查B。下一操作可以分别为细化A与检查B；仅看最后片段无法作出这种区分。

\section*{第\ref{ch:identifiability}章}
1．三个类分别贡献$0$、$0.3/2=0.15$与$0.3/3=0.1$，总误差为$0.25$。

2．添加实验只会为等价关系增加等同性约束，所以每个新类都是旧类的子集。可以忽略新增实验继续旧策略，故最优风险不增；也可直接比较类内最大先验质量。

\section*{第\ref{ch:coverage}章}
1．第一实验权重$w$时，两类覆盖分别为$aw$与$b(1-w)$。最优点令二者相等，故$w=b/(a+b)$，值为$ab/(a+b)$。

2．错误输出必有一个目标不同的情境在似然上击败真实情境，因此情境比较足以保证目标正确。对许多具有相同目标的情境逐个作并合界可能多付出$\ln m$；任务专用读出与表示有机会改进它。

\section*{第\ref{ch:lower}章}
1．使用附录中的两个偏导符号，将$p$降至$1-\delta$、将$q$升至$\delta$都不增大二元相对熵。

2．停止透露第一次观察属于预定停止集合。补空操作保留实际停止位置，完整记录分布仍包含这项信息；若删除位置，只比较存在的输出，便可能遗漏它。

\section*{第\ref{ch:operational}章}
1．例如$M$随规模指数增长，$C_M$也增长。带有$C_M$的下界可能变弱，不能再视为统一常数因子的匹配。

2．以$\lceil\ln(m/\delta)/u\rceil$次调用运行公开分配，可以保证最坏错误不超过$\delta$；另行报告分配取得、控制和通信费用。未知最优值不妨碍可执行上界，但不能据此声称最优。

\section*{第\ref{ch:calibration}章}
1．$(1,0,1)$是两行之和，可辨别。两行均与$(1,1,1)$正交，但目标$(1,1,1)$与该向量内积为1，因此后者不可辨别。

2．给目标涉及的$k$个基向量等权，任一目标冲突差至少命中其中一个坐标，得到$1/k$。反方等权选择这些单坐标差，任何实验的平均命中至多$1/k$，由对偶得最优。

\section*{第\ref{ch:communication}章}
1．所有情境的接收分布均为该符号的点质量，所以所有区别量为零，任意分配覆盖率都为零。

2．公平目标已被A知道，A仅在目标为1时发送固定标记。B看到沉默后应以概率1判断0；如果把沉默当作没有发生任何条件化，它会错误地保留先验$1/2$。

\section*{第\ref{ch:representation}章}
1．总分数误差不超过$E=\sum_je_j$，风险至多$\min\{1,me^E(1-u)^n\}$。若要求不超过$\delta$，一个充分条件为$E\le-n\ln(1-u)-\ln(m/\delta)$，右侧需非负。

2．可逆变换不合并情境似然区别，但逆变换可能很昂贵。例如将一个容易直接读取的答案嵌入需要大量计算才能解码的可逆表示，信息不丢失而受限读出变难。

\section*{第\ref{ch:interpretation}章}
1．成员独立产生噪声，但都观察$Y\oplus B$，共享持续偏差$B$。给定$(Y,B)$的独立性成立；相对于目标$Y$的共同误差并未消除。

2．固定目标冲突对，检查表示是否改变似然比或可辨别类；实现编码与读出并计费；再在全部有限情境或预先指定的实例族上检查风险，而非只复用构造表示时见过的答案。

\section*{第\ref{ch:frontiers}章}
1．第\ref{thm:dimension}条定理的$d=2^r$增长，而$M,\kap$固定。一般主定理的$d$因子变松；补充的信息下界恢复了关于$r$的正确增长率。

2．至少给出合法输入、观察协议、策略允许信息、成本、量词、构造、完整类别下界或反例、可复算数据，以及与最相近结果在假设和结论上的逐项差别。

\chapter{结论、归属与核验}
\label{app:verification}

\section{核心推导的依赖}
\begin{longtable}{p{0.25\textwidth}p{0.34\textwidth}p{0.32\textwidth}}
\toprule
结论 & 核心条件 & 使用的工具\\\midrule
不可辨别误差下限 & 有限静态情境；控制者仅有允许观察 & 记录归纳；有限类识别\\
覆盖率上界 & 独立公开分配；已知概率表 & 乘积亲和度；似然比较\\
全组织下界 & 相同局部策略；严格正核 & KL链式分解；事件数据处理\\
操作性主定理 & 固定有限$d$；统一似然比界 & 上下界；KL与Hellinger比较\\
规模反例 & 随机线性观察；恢复完整情境 & 覆盖计算；Fano；二点下界\\
校准与表示 & 固定二元线性关系；独立调用噪声 & 行空间；覆盖规划；多数集中\\
通信保真 & 相同后处理；完整观察计数 & 柯西不等式；似然因子化\\
停止与模型误差 & 完整证据前缀；明确误差界 & 非负鞅；平方根似然比\\
\bottomrule
\end{longtable}

\section{与已有数学的关系}
统计实验比较提供后处理与决策能力的基本语言；序贯实验和主动假设检验提供选取观察的形式化问题；换测度、Hellinger亲和度、Fano界和鞅界提供概率工具。二元线性观察的可辨别性来自有限域线性代数，多数纠错使用经典有界集中。

本书把任务相对的冲突对、全局覆盖规划和局部KL规划放在同一个有限组织问题中，给出明确常数的连接，再用共同偏差、规模反例和通信变换检验其含义。这些推导的正确性与是否构成新的学术结果分开核对；文献范围记录在随附元数据中。

\section{有限核验怎样使用}
\nolinkurl{scripts/check_theory.py}使用精确分数、有限枚举和浮点容差检查不可辨别误差、覆盖规划、主动记录KL、规模反例、通信压缩、计数充分性与停止概率。它还生成可复算的数据与表格。

\nolinkurl{scripts/check_manuscript.py}检查章节、引文、交叉引用、环境配对和练习数量。构建报告记录XeLaTeX结果与相关页面的目视检查。一般定理由正文证明承担；有限检查用于发现公式、实现与参数口径之间的不一致。

\section{报告数学进展}
可用三个问题检查一项后续改动：它是否扩大了允许的模型而保持结论；是否收紧了整个策略类别的界；是否给出了更经济的可执行组织。只增加一个名称或一个实例时，应说明它承担的是解释还是数学推广。

\backmatter
\bibliographystyle{plainnat}
\bibliography{references}
\end{document}
